\chapter{Bạn Bè Ra Rẽ, Đường Ai Nấy Đi}
\markboth{Bạn Bè Ra Rẽ, Đường Ai Nấy Đi}{Friends Drift, Each Person Has a Road}

\section{Ra Đời Rồi, Bạn Bè Không Còn Đi Cùng Một Tốc Độ}
\begin{truyen}{Ra Đời Rồi, Bạn Bè Không Còn Đi Cùng Một Tốc Độ}{After School, Friends No Longer Move at the Same Speed}
\chuhoa{K}{hi còn đi học, lịch sống của mọi người khá giống nhau.}
\textit{(When people are in school, their schedules are still fairly similar.)}

Ra đời rồi, có người đi làm sớm, có người học tiếp, có người lấy chồng lấy vợ, có người còn đang loay hoay.
\textit{(After that, some start work early, some continue studying, some get married, and some are still searching.)}

Bạn bè vì thế không còn đi cùng một nhịp như trước.
\textit{(So friends no longer move in the same rhythm as before.)}

Nếu không hiểu điều đó, người trẻ rất dễ buồn vì nghĩ rằng tình bạn đang chết đi.
\textit{(If you do not understand this, it is easy to think friendship is dying.)}

\end{truyen}

\begin{baihoc}
Nhiều tình bạn không hỏng. Nó chỉ đổi nhịp. Người trưởng thành phải học chấp nhận nhịp khác nhau mà vẫn giữ lòng tử tế.
\textit{(Many friendships are not broken. They only change pace. Adults must learn to accept different rhythms while keeping kindness.)}
\end{baihoc}

\begin{ghinhoanh}
\textbf{Short English to remember:}

Many friendships are not broken.

They only change pace.

Adults must learn to accept different rhythms while keeping kindness.

\end{ghinhoanh}

\begin{tuhocnhanh}
\textbf{Cau hoi tu hoc / Self-study questions}

1. Van de chinh la gi? \textit{What is the main problem?}

2. Bai hoc thuc te la gi? \textit{What is the practical lesson?}

3. Mot cau tieng Anh em muon nho la cau nao? \textit{Which English sentence do you want to remember?}
\end{tuhocnhanh}
\ngancach

\section{Có Người Chỉ Hợp Đi Cùng Em Một Đoạn}
\begin{truyen}{Có Người Chỉ Hợp Đi Cùng Em Một Đoạn}{Some People Fit Only One Part of the Road}
\chuhoa{C}{ó những người rất hợp với em ở tuổi mười bảy nhưng không còn hợp ở tuổi hai mươi ba.}
\textit{(Some people fit you well at seventeen but no longer fit at twenty-three.)}

Điều đó không nhất thiết ai phản bội ai.
\textit{(That does not always mean betrayal.)}

Con người đổi giá trị, đổi môi trường, đổi nhịp sống, nên tình bạn cũng đổi theo.
\textit{(People change values, environments, and pace, so friendship changes too.)}

Cố giữ một mối quan hệ chỉ vì quá khứ đẹp đôi khi làm hiện tại mệt thêm.
\textit{(Trying to keep a relationship only because the past was beautiful can make the present heavier.)}

\end{truyen}

\begin{baihoc}
Không phải ai từng thân cũng phải ở lại mãi. Có người đến để đi cùng em một quãng, vậy là đủ.
\textit{(Not everyone who was once close must stay forever. Some people come to walk with you for one part, and that is enough.)}
\end{baihoc}

\begin{ghinhoanh}
\textbf{Short English to remember:}

Not everyone who was once close must stay forever.

Some people come to walk with you for one part, and that is enough.

\end{ghinhoanh}

\begin{tuhocnhanh}
\textbf{Cau hoi tu hoc / Self-study questions}

1. Van de chinh la gi? \textit{What is the main problem?}

2. Bai hoc thuc te la gi? \textit{What is the practical lesson?}

3. Mot cau tieng Anh em muon nho la cau nao? \textit{Which English sentence do you want to remember?}
\end{tuhocnhanh}
\ngancach

\section{Thành Công Của Bạn Có Thể Làm Em Chạnh Lòng}
\begin{truyen}{Thành Công Của Bạn Có Thể Làm Em Chạnh Lòng}{A Friend’s Success May Sting You}
\chuhoa{R}{a đời rồi, chênh lệch giữa bạn bè hiện ra rất nhanh.}
\textit{(After school, the gap between friends can appear very quickly.)}

Người này có việc tốt, người kia thất nghiệp.
\textit{(One person gets a good job, another is unemployed.)}

Người này mua xe, người kia còn nợ trọ.
\textit{(One buys a motorbike, another still owes rent.)}

Chạnh lòng trước thành công của bạn là rất người.
\textit{(Feeling a sting before a friend’s success is very human.)}

Điều quan trọng là đừng để cơn chạnh lòng đó làm em ghét điều tốt của người khác.
\textit{(What matters is not letting that sting turn into hatred for someone else’s good fortune.)}

\end{truyen}

\begin{baihoc}
So sánh là điều khó tránh. Nhưng nếu giữ được lòng sạch trước thành công của bạn, em sẽ đỡ cay độc và đỡ khổ hơn.
\textit{(Comparison is hard to avoid. But if you keep your heart clean before a friend’s success, you will become less bitter and less miserable.)}
\end{baihoc}

\begin{ghinhoanh}
\textbf{Short English to remember:}

Comparison is hard to avoid.

But if you keep your heart clean before a friend’s success, you will become less bitter and less miserable.

\end{ghinhoanh}

\begin{tuhocnhanh}
\textbf{Cau hoi tu hoc / Self-study questions}

1. Van de chinh la gi? \textit{What is the main problem?}

2. Bai hoc thuc te la gi? \textit{What is the practical lesson?}

3. Mot cau tieng Anh em muon nho la cau nao? \textit{Which English sentence do you want to remember?}
\end{tuhocnhanh}
\ngancach

\section{Tin Nhắn Thưa Dần Không Luôn Là Hết Thân}
\begin{truyen}{Tin Nhắn Thưa Dần Không Luôn Là Hết Thân}{Fewer Messages Do Not Always Mean Less Love}
\chuhoa{L}{ớn lên rồi, có những người thương nhau thật nhưng nói chuyện ít đi.}
\textit{(As people grow older, some truly care for each other yet speak less often.)}

Không phải vì họ hết quý nhau.
\textit{(It is not always because they care less.)}

Mà vì đời sống bắt đầu đầy trách nhiệm khác.
\textit{(It is because life fills with other responsibilities.)}

Người trẻ nếu quá bám vào tần suất liên lạc sẽ dễ hiểu lầm nhiều mối quan hệ tốt.
\textit{(If young people cling too much to message frequency, they may misread many good relationships.)}

Có những tình bạn trưởng thành hơn khi nó bớt ồn mà vẫn còn thật.
\textit{(Some friendships become more mature when they grow less noisy but remain true.)}

\end{truyen}

\begin{baihoc}
Tình bạn lớn lên không phải lúc nào cũng ồn ào. Có khi nó chỉ trở nên lặng hơn nhưng chắc hơn.
\textit{(Friendship does not always grow louder as it matures. Sometimes it grows quieter but stronger.)}
\end{baihoc}

\begin{ghinhoanh}
\textbf{Short English to remember:}

Friendship does not always grow louder as it matures.

Sometimes it grows quieter but stronger.

\end{ghinhoanh}

\begin{tuhocnhanh}
\textbf{Cau hoi tu hoc / Self-study questions}

1. Van de chinh la gi? \textit{What is the main problem?}

2. Bai hoc thuc te la gi? \textit{What is the practical lesson?}

3. Mot cau tieng Anh em muon nho la cau nao? \textit{Which English sentence do you want to remember?}
\end{tuhocnhanh}
\ngancach

\section{Cũng Có Lúc Phải Chia Tay Một Tình Bạn}
\begin{truyen}{Cũng Có Lúc Phải Chia Tay Một Tình Bạn}{Sometimes You Must End a Friendship}
\chuhoa{C}{ó những người bạn chỉ hợp khi cả hai còn vô tư.}
\textit{(Some friendships work only while both people remain carefree.)}

Khi đời bắt đầu khó hơn, những lệch giá trị, lệch cách sống, và lệch đạo đức hiện ra rõ hơn.
\textit{(When life becomes harder, differences in values, lifestyle, and ethics become clearer.)}

Không phải mối quan hệ nào cũng cần cứu bằng mọi giá.
\textit{(Not every relationship should be saved at all costs.)}

Có những cuộc chia tay tử tế giúp cả hai bớt làm hỏng nhau hơn.
\textit{(Some decent endings keep both people from damaging each other further.)}

\end{truyen}

\begin{baihoc}
Người trưởng thành phải biết có những tình bạn nên được buông xuống đàng hoàng, không phải kéo lê mãi trong khó chịu.
\textit{(A mature person must know that some friendships should be laid down with dignity, not dragged on in resentment.)}
\end{baihoc}

\begin{ghinhoanh}
\textbf{Short English to remember:}

A mature person must know that some friendships should be laid down with dignity, not dragged on in resentment.

\end{ghinhoanh}

\begin{tuhocnhanh}
\textbf{Cau hoi tu hoc / Self-study questions}

1. Van de chinh la gi? \textit{What is the main problem?}

2. Bai hoc thuc te la gi? \textit{What is the practical lesson?}

3. Mot cau tieng Anh em muon nho la cau nao? \textit{Which English sentence do you want to remember?}
\end{tuhocnhanh}
\ngancach

\section{Bạn Chỉ Tìm Đến Khi Cần}
\begin{truyen}{Bạn Chỉ Tìm Đến Khi Cần}{A Friend Who Appears Only When Needing Something}
\chuhoa{C}{ó những tình bạn tưởng thân cho đến khi em nhận ra người kia chỉ tìm tới mỗi lúc cần nhờ vả.}
\textit{(Some friendships feel close until you realize the other person appears only when they need something.)}

Em từng tưởng cứ quen lâu là tình nghĩa tự sâu.
\textit{(You once thought time alone automatically deepened the bond.)}

Nhưng thời gian không tự làm nên sự tử tế.
\textit{(But time does not automatically create goodness.)}

Nó chỉ làm rõ người ta thật sự coi em là gì.
\textit{(It only reveals what you truly mean to someone.)}

Nhìn ra điều đó buồn, nhưng nó giúp em bớt đặt lòng tốt sai chỗ.
\textit{(Seeing that is sad, but it helps you place your kindness better.)}

\end{truyen}

\begin{baihoc}
Không phải ai biết số em cũng là bạn của em.
\textit{(Not everyone who has your number is your friend.)}
\end{baihoc}

\begin{ghinhoanh}
\textbf{Short English to remember:}

Not everyone who has your number is your friend.

\end{ghinhoanh}

\begin{tuhocnhanh}
\textbf{Cau hoi tu hoc / Self-study questions}

1. Van de chinh la gi? \textit{What is the main problem?}

2. Bai hoc thuc te la gi? \textit{What is the practical lesson?}

3. Mot cau tieng Anh em muon nho la cau nao? \textit{Which English sentence do you want to remember?}
\end{tuhocnhanh}
\ngancach

\section{Tình Bạn Lớn Lên Cũng Cần Công Chăm}
\begin{truyen}{Tình Bạn Lớn Lên Cũng Cần Công Chăm}{Adult Friendship Also Needs Care}
\chuhoa{N}{hiều người buồn vì nghĩ bạn cũ không chủ động nghĩa là tình cảm đã hết.}
\textit{(Many feel hurt and think that old friends have stopped caring because they no longer initiate contact.)}

Họ tưởng tình bạn thật sẽ tự sống mãi không cần vun.
\textit{(They think true friendship should survive forever without effort.)}

Nhưng bạn bè trưởng thành cũng cần sự chủ động, lời hỏi han, và vài lần đi về phía nhau.
\textit{(But adult friendship also needs initiative, small check-ins, and both people walking toward each other sometimes.)}

Giữ một tình bạn tốt đôi khi cần cố ý, không chỉ dựa vào thói quen cũ.
\textit{(Keeping a good friendship sometimes requires intention, not only old habit.)}

\end{truyen}

\begin{baihoc}
Tình bạn bền không chỉ ở ký ức đẹp. Nó ở sự chăm lại khi đời đã bận.
\textit{(Stable friendship lives not only in memory, but in renewed care after life becomes busy.)}
\end{baihoc}

\begin{ghinhoanh}
\textbf{Short English to remember:}

Stable friendship lives not only in memory, but in renewed care after life becomes busy.

\end{ghinhoanh}

\begin{tuhocnhanh}
\textbf{Cau hoi tu hoc / Self-study questions}

1. Van de chinh la gi? \textit{What is the main problem?}

2. Bai hoc thuc te la gi? \textit{What is the practical lesson?}

3. Mot cau tieng Anh em muon nho la cau nao? \textit{Which English sentence do you want to remember?}
\end{tuhocnhanh}
\ngancach

\section{Thu Nhập Khác, Vòng Tròn Khác}
\begin{truyen}{Thu Nhập Khác, Vòng Tròn Khác}{Different Incomes Create Different Circles}
\chuhoa{K}{hi thu nhập mỗi người khác nhau rõ hơn, nhiều buổi hẹn bắt đầu trở nên khó xử.}
\textit{(When incomes become more different, many meetups begin to feel awkward.)}

Nhiều người tưởng nếu thân thật thì tiền bạc sẽ không bao giờ chen vào.
\textit{(Many think real friendship should never be touched by money differences.)}

Nhưng điều kiện sống ảnh hưởng đến nơi người ta đi, cách người ta tiêu, và cả cảm giác thoải mái khi gặp nhau.
\textit{(But living conditions affect where people go, how they spend, and even how relaxed they feel together.)}

Hiểu điều đó giúp em bớt tự ái và bớt ép người khác phải giữ nhịp như cũ.
\textit{(Understanding that helps you feel less offended and stop forcing old rhythms.)}

\end{truyen}

\begin{baihoc}
Bạn bè lớn lên sẽ có lúc đi qua những khoảng chênh rất thật.
\textit{(Friends growing older must sometimes cross very real gaps.)}
\end{baihoc}

\begin{ghinhoanh}
\textbf{Short English to remember:}

Friends growing older must sometimes cross very real gaps.

\end{ghinhoanh}

\begin{tuhocnhanh}
\textbf{Cau hoi tu hoc / Self-study questions}

1. Van de chinh la gi? \textit{What is the main problem?}

2. Bai hoc thuc te la gi? \textit{What is the practical lesson?}

3. Mot cau tieng Anh em muon nho la cau nao? \textit{Which English sentence do you want to remember?}
\end{tuhocnhanh}
\ngancach

\section{Lời Hứa Cũ Phải Sống Trong Hiện Tại}
\begin{truyen}{Lời Hứa Cũ Phải Sống Trong Hiện Tại}{Old Promises Must Still Survive the Present}
\chuhoa{N}{gày trước ai cũng từng hứa sẽ không bỏ nhau, nhưng rồi cuộc sống kéo từng người vào một hướng khác.}
\textit{(In the past everyone promised never to leave, but life later pulled each person in a different direction.)}

Em tưởng lời hứa cũ đủ mạnh để giữ mọi thứ nguyên như trước.
\textit{(You think old promises are strong enough to keep things unchanged.)}

Nhưng lời hứa chỉ sống được khi hiện tại còn có hành động nuôi nó.
\textit{(But a promise lives only when present actions still feed it.)}

Không phải lời hứa nào cũng bị phản bội.
\textit{(Not every promise is betrayed.)}

Có khi chỉ là nó không đủ sức đi qua một đời sống quá khác.
\textit{(Sometimes it simply cannot survive a very different life.)}

\end{truyen}

\begin{baihoc}
Người trưởng thành biết trân trọng lời hứa đẹp mà vẫn nhìn thẳng vào thực tại.
\textit{(A mature person respects beautiful promises while still facing reality.)}
\end{baihoc}

\begin{ghinhoanh}
\textbf{Short English to remember:}

A mature person respects beautiful promises while still facing reality.

\end{ghinhoanh}

\begin{tuhocnhanh}
\textbf{Cau hoi tu hoc / Self-study questions}

1. Van de chinh la gi? \textit{What is the main problem?}

2. Bai hoc thuc te la gi? \textit{What is the practical lesson?}

3. Mot cau tieng Anh em muon nho la cau nao? \textit{Which English sentence do you want to remember?}
\end{tuhocnhanh}
\ngancach

\section{Giữ Lại Sự Tôn Trọng Sau Khi Tan Ra}
\begin{truyen}{Giữ Lại Sự Tôn Trọng Sau Khi Tan Ra}{Keeping Respect After Drifting Apart}
\chuhoa{C}{ó những tình bạn không còn gần nữa nhưng vẫn không cần biến thành cay cú hay mỉa mai.}
\textit{(Some friendships are no longer close, yet they do not need to turn bitter or mocking.)}

Nhiều người tưởng đã không đi cùng được thì phải kết thúc bằng tổn thương.
\textit{(Many people think that if a friendship cannot continue closely, it must end in hurt.)}

Thật ra có những mối quan hệ đẹp nhất ở chỗ nó biết rời xa một cách tử tế.
\textit{(In truth, some relationships are most beautiful in the way they know how to part with decency.)}

Giữ sự tôn trọng sau khi không còn thân là một dấu hiệu rất trưởng thành.
\textit{(Keeping respect after distance is a very mature sign.)}

\end{truyen}

\begin{baihoc}
Không phải tình bạn nào cũng cần ở lại. Nhưng tình người thì vẫn có thể được giữ lại.
\textit{(Not every friendship must remain close. Human respect can still remain.)}
\end{baihoc}

\begin{ghinhoanh}
\textbf{Short English to remember:}

Not every friendship must remain close.

Human respect can still remain.

\end{ghinhoanh}

\begin{tuhocnhanh}
\textbf{Cau hoi tu hoc / Self-study questions}

1. Van de chinh la gi? \textit{What is the main problem?}

2. Bai hoc thuc te la gi? \textit{What is the practical lesson?}

3. Mot cau tieng Anh em muon nho la cau nao? \textit{Which English sentence do you want to remember?}
\end{tuhocnhanh}
